% Detailed PNM slides retained from the main presentation.
\begin{frame}{Pore-network geometry and throat diffusion}
\setlength{\abovedisplayskip}{5pt}\setlength{\belowdisplayskip}{5pt}
  \begin{columns}[T,onlytextwidth]
    \column{0.46\textwidth}
      \hi{Network construction}
      \footnotesize
      \begin{itemize}\setlength{\itemsep}{1pt}\setlength{\parskip}{0pt}
        \item four bead centres define a Delaunay tetrahedron; its valid circumcentre represents a pore node
        \item tetrahedra sharing a triangular face are joined by a throat
        \item topology comes from particle coordinates, not voxels
      \end{itemize}
      \vspace{0.12cm}
      \begin{itemize}\setlength{\itemsep}{1pt}\setlength{\parskip}{0pt}
        \item $2^{22}$ Sobol points sample the solid and gas space.
        \item Gas-point counts allocate the calculated void volume to pores.
      \end{itemize}
    \column{0.50\textwidth}
      \hi{Diffusion through throat $i$--$j$}
      \[
        G_{ij}=D_m\frac{A_{ij}}{L_{ij}},\qquad
        \dot n_{ij}=G_{ij}(c_j-c_i).
      \]
      \footnotesize
      $D_m=1.5\times10^{-5}$ m$^2$ s$^{-1}$: molecular CO$_2$ diffusivity in the gas.
      \par\smallskip
      $L_{ij}=\lVert\mathbf{x}_j-\mathbf{x}_i\rVert$ is the pore-centre distance.
      The irregular opening between three beads is approximated by an equivalent
      circular area
      \[
        A_{ij}=\pi r_{ij}^{2},
      \]
      where $r_{ij}$ is the clearance of the shared Delaunay face.
  \end{columns}
  \vspace{0.16cm}
  \result{The circular throat area is a geometric modelling assumption; $G_{ij}$ has units of m$^3$ s$^{-1}$.}
\end{frame}

\begin{frame}{The pore network solves diffusion and adsorption together}

\begin{columns}[T,onlytextwidth]
\column{0.56\textwidth}\small
\hi{Gas balance for each pore $i$}
\[
V_i\frac{dc_i}{dt}=\sum_{j\in\mathcal N(i)}G_{ij}(c_j-c_i)
+G_{iR}(c_R-c_i)-S_i.
\]
\begin{itemize}\setlength{\itemsep}{2pt}
\item $V_i$: gas volume, $c_i$: gas concentration
\item $G_{ij}$: throat conductance [m$^3$ s$^{-1}$]
\item $G_{iR}$: connection to reservoir at $c_R$, zero away from the inlet
\item $S_i$: net gas transfer to adjacent beads [mol s$^{-1}$]
\end{itemize}
\column{0.40\textwidth}\small
\hi{Uptake for each bead $k$}
\[
\frac{dq_k}{dt}=k_{\rm LDF}(q_k^*-q_k).
\]
\begin{itemize}\setlength{\itemsep}{2pt}
\item $q_k$: current loading [mol kg$^{-1}$]
\item $k_{\rm LDF}$ [s$^{-1}$]: linear-driving-force rate coefficient
\item $1/k_{\rm LDF}$: relaxation time at fixed gas concentration
\end{itemize}
\end{columns}
\vspace{0.15cm}
\result{One coupled system: one concentration per pore and one loading per bead.}

\end{frame}

\begin{frame}{Local equilibrium sets the adsorption target}
\begin{columns}[T,onlytextwidth]
\column{0.48\textwidth}\hi{Local equilibrium loading}\small
\[
p_k=c_k\:R\:T,\qquad q_k^*=\frac{q_s\:b\:p_k}{1+b\:p_k}.
\]
\begin{itemize}\setlength{\itemsep}{3pt}
\item $c_k$: average gas concentration around bead $k$
\item $p_k$: local CO$_2$ partial pressure [Pa]
\item $q_k^*$: equilibrium target [mol kg$^{-1}$]
\end{itemize}
\column{0.48\textwidth}\hi{Fixed material parameters}\small
\begin{itemize}\setlength{\itemsep}{3pt}
\item $q_s$: maximum CO$_2$ loading [mol kg$^{-1}$]
\item $b$: Langmuir affinity [Pa$^{-1}$], so $b\:p_k$ is dimensionless
\item $k_{\rm LDF}$: kinetic coefficient [s$^{-1}$]
\end{itemize}
\medskip
\hi{Actual loading evolves in time}
\[
\frac{dq_k}{dt}=k_{\rm LDF}(q_k^*-q_k).
\]
\end{columns}
\vspace{0.25cm}
\result{Local gas concentrations set the target. Kinetics determine how quickly beads approach it.}
\end{frame}

\begin{frame}{The pore network delivers the reference uptake curve}

\begin{columns}[T,onlytextwidth]
\column{0.47\textwidth}\hi{Resolved transient outputs}\small
\begin{itemize}
\item Gas concentration $c_i(t)$ in every pore
\item Adsorbed loading $q_k(t)$ of every bead
\item Different bead loadings reveal incomplete utilization
\end{itemize}
\medskip
The detailed model retains the actual pore connections and bead associations.
\column{0.49\textwidth}\hi{Bed-average uptake}\small
\[
\overline q_{\rm PNM}(t)=\frac{\sum_k m_k q_k(t)}{\sum_k m_k}.
\]
For identical beads:
\[
\overline q_{\rm PNM}(t)=\frac{1}{N_b}\sum_{k=1}^{N_b}q_k(t).
\]
This curve is the reference for fitting the reduced 1D model.
\end{columns}
\vspace{0.25cm}
\result{Next: can a homogeneous 1D bed reproduce this complete uptake curve?}

\end{frame}

